\documentclass[11pt]{amsart}
\usepackage{amsmath,amssymb,amsthm,mathtools}
\newtheorem{theorem}{Theorem}[section]
\newtheorem{proposition}[theorem]{Proposition}
\newtheorem{definition}[theorem]{Definition}
\newtheorem{remark}[theorem]{Remark}
\title{Affine Normal-Cone Flatification and Tor--Valabrega Transform Factorization}
\author{Maximum-Likelihood Candidate Chapter}
\date{4 August 2026}
\begin{document}
\maketitle

\begin{abstract}
We isolate a finite obstruction to the interchange of ambient strict transform
and normal associated grading along a regular carrier.  On every affine chart
of a centre-enriched trace blowup, the comparison factors canonically through
the controlled-transform filtration.  The first map is an epimorphism whose
kernel is a Tor/base-change defect.  The second map may have both kernel and
cokernel; these are Valabrega-type intersection and saturation-generation
defects.  Vanishing of all three defects gives the transform-interchange
isomorphism.  The complete associated-graded and defect portfolios are finite
modules on affine normal and mixed-Rees schemes, so general affine
flatification applies without a projective Rees--Serre tail.  The ambient
realization and positive-characteristic resolution theorem remain open.
\end{abstract}

\section{The affine chart and relative saturation}
Let $A$ be Noetherian, let $I\subset K\subset A$ be ideals, let $M$ be a
finite $A$-module, and choose $q\in K$.  Put
\[
 B=A[K/q],\qquad N=M\otimes_A B,
 \qquad Q_q(P)=P/H^0_{(q)}(P).
\]
For $P\subset N$, define
\[
 \operatorname{Sat}_q(P;N)=
 \{x\in N:\ q^n x\in P\text{ for some }n\geq0\}.
\]
Then
\[
 H^0_{(q)}(N/P)=\operatorname{Sat}_q(P;N)/P,
 \qquad
 Q_q(N/P)\cong N/\operatorname{Sat}_q(P;N).
\]

Let $\overline N=Q_q(N)$.  The exceptional parameter $q$ is regular on
$\overline N$.  Let $J$ be the controlled transform of $I$ on the $q$-chart,
generated by the elements $i/q$ for $i\in I$, and let
\[
 L=\operatorname{Sat}_q(J;B)
\]
be the ideal of the strict transform of the carrier.

\section{The controlled-transform core}
For $a\geq0$, let $O_a$ be the Cartier-twisted strict transform of the old
normal graded piece
\[
 (I^aM/I^{a+1}M)\otimes_A B.
\]
The controlled filtration on $\overline N$ gives
\[
 C_a=J^a\overline N/J^{a+1}\overline N,
 \qquad
 A_a=L^a\overline N/L^{a+1}\overline N.
\]
There is a canonical factorization
\[
 O_a\xrightarrow{\alpha_a}C_a\xrightarrow{\beta_a}A_a.
\]
The first map is surjective.  Its kernel is the filtration-base-change or Tor
defect.  The second map has
\[
 \ker(\beta_a)=
 \frac{J^a\overline N\cap L^{a+1}\overline N}
      {J^{a+1}\overline N}
\]
and
\[
 \operatorname{coker}(\beta_a)=
 \frac{L^a\overline N}
      {J^a\overline N+L^{a+1}\overline N}.
\]
These are the Valabrega intersection defect and the saturation-generation
defect.

\begin{remark}[Why two epimorphisms are false]
Take $B=k[q,y]$, $J=(qy)$, and $L=\operatorname{Sat}_q(J)=(y)$.  The map
$J/J^2\to L/L^2$ has image $q(L/L^2)$ and is not surjective.  Thus the first
X038 cospan candidate is false before saturation compatibility has been
proved.  The factorization above is the corrected primitive.
\end{remark}

\begin{proposition}[Transform-factorization compiler]
If $\ker\alpha_a=0$, $\ker\beta_a=0$, and
$\operatorname{coker}\beta_a=0$, then the composite induces a canonical
isomorphism $O_a\cong A_a$.
\end{proposition}

\begin{proof}
The first map is then a bijection and the second map is a bijection.  Compose
the corresponding linear equivalences.
\end{proof}

\section{Affine normal and mixed-Rees packets}
The algebra
\[
 \operatorname{gr}_I(A)=\bigoplus_{a\geq0}I^a/I^{a+1}
\]
is a finite-type Noetherian algebra over $A/I$, and
$\operatorname{gr}_I(M)$ is finite over it.  Thus the complete passive
associated-graded portfolio is one coherent module on the affine normal cone
\[
 \operatorname{Spec}_{A/I}\operatorname{gr}_I(A).
\]

Likewise,
\[
 \mathcal R_{I,K}(A)=
 \bigoplus_{a,b\geq0}I^aK^bU^aV^b,
 \qquad
 \mathcal R_{I,K}(M)=
 \bigoplus_{a,b\geq0}I^aK^bM U^aV^b
\]
are respectively a finite-type Noetherian algebra and a finite module.  Kernels,
cokernels, power-torsion modules, and stabilized colon quotients arising from
homogeneous maps are finite.  Therefore all degrees of the Tor--Valabrega
factorization form one finite affine coherent packet.  No projective tail or
Hilbert-polynomial cutoff is required for flatifier existence.

\section{Affine flat-kill legalization}
Let $C$ be a regular carrier and let $\mathcal T$ be the finite direct sum of
the affine normal-cone passive portfolio, the Tor kernel, the two Valabrega
defects, log-conormal defects, and source-Cartier torsion.  On the open locus
where the proposed centre is already legal, the defect summands vanish and the
passive summands are flat.  A $U$-admissible flatifying blowup makes the complete
strict transform finitely presented and flat.  A finitely presented flat
defect module which vanishes on a componentwise dense open is zero.

The flatifier centre need not be regular.  Its source ideal is principalized on
$C$, and every regular trace subcentre is recursively legalized before ambient
blowup.  Since each such subcentre has strictly smaller carrier dimension, this
is a well-founded nested recursion.

\begin{theorem}[Maximum-likelihood decisive theorem]
Assume strong relative principalization and the same legalization theorem for
regular carriers of dimension strictly less than $d$.  Let $C$ be a regular
$d$-dimensional carrier in a smooth ambient scheme, equipped with a finite
source-labelled active, passive, logarithmic, and contact portfolio.  Then the
complete affine normal-cone and Tor--Valabrega factorization packet admits a
finite centre-enriched flat-kill word in ordinary ambient blowups with regular,
recursively legalized trace centres.  After the word, all Tor and Valabrega
defects vanish, the ambient strict transform and the new normal associated
graded module agree on every chart and overlap, and the transformed carrier is
jointly legal.  All source, owner, boundary, and debt identities reenter without
reset.
\end{theorem}

\begin{remark}
The preceding theorem is a target statement.  It has not been proved.  In
particular, the scheme-level factorization maps, overlap gluing, regular ambient
realization of an arbitrary flatifier, and hereditary no-reset remain open.
\end{remark}

\section{Publication architecture}
The current maximum-likelihood series consists of six papers totaling 552
dense Annals/AMS-equivalent pages.  The present chapter belongs to Paper III,
\emph{Affine Normal-Cone Flatification, Tor--Valabrega Factorization, and
Hereditary Reentry}.

\end{document}
